\documentclass[12pt,a4paper]{article}
\usepackage[T1]{fontenc}
\usepackage[spanish]{babel}
\usepackage[margin=2.5cm]{geometry}
\usepackage{graphicx,float,booktabs,tabularx,csquotes,hyperref}
\usepackage[backend=biber,style=ieee]{biblatex}
\IfFileExists{referencias.bib}{\addbibresource{referencias.bib}}{}
\newcommand{\evidencia}[3]{\begin{figure}[H]\centering\IfFileExists{#1}{\includegraphics[width=\linewidth,height=.65\textheight,keepaspectratio]{#1}}{\fbox{Imagen pendiente de carga}}\caption{#2}\label{#3}\end{figure}}
\begin{document}
\begin{titlepage}\centering
{\Large Universidad Técnica Estatal de Quevedo\par}
Facultad de Ciencias de la Computación (FCC)\par
Carrera de Ingeniería de Software\par
\vfill
Aplicaciones Distribuidas — TA-IND-01\par
\vspace{1cm}
{\LARGE\bfseries Análisis crítico de Apache Cassandra\par}
Replicación, consistencia y disponibilidad de lectura\par
\vfill
Kevin Leonel García Bazurto\par
Paralelo A\par
13 de noviembre de 2026\par
\href{https://github.com/Yukisan-uwo/TA-IND-01-Informe-Academico-Individual-Analisis-Critico-de-Sistemas-Distribuidos-Grupo}{Repositorio en GitHub}
\end{titlepage}
\begin{abstract}
% Redactar al final: 150–200 palabras.
\end{abstract}
\section{Introducción}
% Problema, usuarios, escala y justificación con fuentes.
\section{Arquitectura del sistema}
% Componentes, comunicación, diagrama propio y evolución desde el Anexo A.
\section{Análisis con los conceptos de la Unidad 1}
\subsection{Transparencias de distribución}
\subsection{CAP y PACELC}
\subsection{Modelos de comunicación}
\subsection{Tolerancia a fallos}
\section{Experimentación propia}
\subsection{Objetivo y configuración}
La prueba evaluó si Cassandra podía devolver el mismo registro con \texttt{ONE} y \texttt{QUORUM} al reducir el número de réplicas activas. El entorno usó Cassandra 5.0.9 en Docker, tres nodos en \texttt{dc1} y un factor de replicación de tres.
\evidencia{Evidencias/01-cluster-tres-nodos.jpeg}{Estado inicial del clúster con tres nodos activos.}{fig:cluster}
\subsection{Procedimiento y resultados}
El registro con \texttt{id = 1} se insertó con \texttt{ALL}, que exigía la confirmación de las tres réplicas. Se mantuvo activo \texttt{taind-cas1} como coordinador de la consulta. Los comandos y las salidas se documentan en \path{Experimento/prueba-cassandra.md}.




Con tres réplicas activas, ambas lecturas devolvieron el registro.
\evidencia{Evidencias/02-lecturas-tres-replicas.png}{Lecturas con tres réplicas activas.}{fig:tres}




Después de detener \texttt{taind-cas3}, se repitieron ambas lecturas con dos réplicas activas (Figura~\ref{fig:dos}).
\evidencia{Evidencias/03-lecturas-dos-replicas.png}{Lecturas tras detener el tercer nodo.}{fig:dos}




Al detener también \texttt{taind-cas2}, \texttt{ONE} devolvió el registro y \texttt{QUORUM} produjo un error \texttt{Unavailable}. Cassandra indicó dos réplicas requeridas y una disponible.
\evidencia{Evidencias/04-lecturas-una-replica.png}{Lecturas con una réplica activa.}{fig:una}
\subsection{Interpretación}
En este escenario, \texttt{ONE} pudo completar la lectura con la réplica disponible. \texttt{QUORUM} requería dos respuestas de las tres réplicas configuradas. El resultado muestra cómo el nivel de consistencia afectó la disponibilidad de la lectura durante la detención controlada de nodos.
\section{Ejemplos del entorno del estudiante}
% Dos observaciones propias y su relación con los conceptos de la unidad.
\section{Análisis comparativo}
% Confirmar el comparador y discutir al menos cinco criterios con fuentes.
\section{Postura crítica}
\section{Conclusiones}
\section{Declaración de uso de inteligencia artificial}
% Identificar herramienta, versión verificable, usos y comprobaciones.
\IfFileExists{referencias.bib}{\printbibliography[title={Referencias}]}{}
\clearpage
\section*{Anexos}
\appendix
\section{Esquema inicial a mano}
\evidencia{esquema-inicial.jpeg}{Bosquejo inicial del clúster Cassandra.}{fig:esquema}
\section{Evidencias de lectura de las fuentes}
% Tres capturas anotadas con su página y cita correspondiente.
\end{document}